\SourceLeaf{61}
$\gamma$) Sind diese Koeffizienten bekannt, so lassen sich die endlichen Ausdrücke für die übrigen leicht mit dem zuvor gefundenen Rekursionsschema herleiten. Dieses lautete
\[
A^{(h+2)} = \frac{n(\nu-h)A^{(h)} - 2(h+1)A^{(h+1)}}{n(\nu+h+2)}
\]
also hier wegen $\nu=-1$
\[
A^{(h+2)} = \frac{-n(h+1)A^{(h)} - 2(h+1)A^{(h+1)}}{n(h+1)}
\]
oder
\[
A^{(h+2)} = -\Bigl(A^{(h)} + \frac{2}{n}A^{(h+1)}\Bigr).
\]
Führt man die Rechnungen nach diesem Schema aus, so stößt man auf ein sehr einfaches Gesetz\HistoricalNote{*)}{*)}{Der Ausdruck für dieses Gesetz bei \textsc{Euler}, a. a. O., \S\,275 ist falsch, vielleicht aufgrund eines Setzfehlers.}, nämlich
\[
A^{(h)} = (-1)^h\,\frac{2}{\BareRoot{(1-nn)}}\left(\frac{1-\BareRoot{(1-nn)}}{n}\right)^{h}.
\]
Um es zu beweisen, nehmen wir an, es gelte für die Koeffizienten bis zum $h^{\text{tum}}$ Koeffizienten. Nach dem Rekursionsschema erhält man
\[
\begin{aligned}
A^{(h+1)} & = (-1)^h\,\frac{2}{\BareRoot{(1-nn)}}\left(\frac{1-\BareRoot{(1-nn)}}{n}\right)^{h-1}\\&\quad + (-1)^{h+1}\,\frac{2}{n}\,\frac{2}{\BareRoot{(1-nn)}}\left(\frac{1-\BareRoot{(1-nn)}}{n}\right)^{h}
\end{aligned}
\]
oder, wenn man den gemeinsamen Faktor ausklammert,
\[
(-1)^h\,\frac{2}{\BareRoot{(1-nn)}}\left(\frac{1-\BareRoot{(1-nn)}}{n}\right)^{h-1}\left\{1-\frac{2}{n}\left(\frac{1-\BareRoot{(1-nn)}}{n}\right)\right\}.
\]
Der in geschweiften Klammern stehende Faktor ist
\[
\begin{aligned}
&= \frac{nn-2+2\BareRoot{(1-nn)}}{nn}\\&= -\frac{(1-nn)-2\BareRoot{(1-nn)}+1}{nn}\\&= -\left(\frac{1-\BareRoot{(1-nn)}}{n}\right)^{2},
\end{aligned}
\]
Setzt man ihn ein, so ergibt sich
\[
\begin{aligned}
A^{(h+1)} & = (-1)^h\,\frac{2}{\BareRoot{(1-nn)}}\left(\frac{1-\BareRoot{(1-nn)}}{n}\right)^{h-1}\\&\quad .  -\left(\frac{1-\BareRoot{(1-nn)}}{n}\right)^{2}\\&= (-1)^{h+1} . \frac{2}{\BareRoot{(1-nn)}}\left(\frac{1-\BareRoot{(1-nn)}}{n}\right)^{h+1},
\end{aligned}
\]
woraus die allgemeine Gültigkeit des Gesetzes ersichtlich ist.

